% ===========================================================================
\section{Introduction}
% ===========================================================================

\Acf{amt} has been proposed as a way to shape model behavior before post-training. Recent work suggests that training base models on alignment-relevant, pretraining-style documents can implant behavior-relevant knowledge
\citep{wang2025modifying,mcdougall2026synthetic}, steer downstream behavior \citep{tice2026alignment}, and shape how models generalize from underspecified finetuning data \citep{li2026modelspec}. Some of these effects persist through subsequent \ac{sft} and \ac{rl} \citep{kutasov2026teaching}. Much of the motivation for pursuing alignment midtraining is that it could compensate for limitations of post-training. In practice, post-training demonstrations may contain conflicting objectives or may fail to cover every example of the behavior we want a model to learn. Midtraining aims to instill the underlying principles, motivations, or character that generate those behaviors. The hope is that models will then apply these principles in out-of-distribution settings, even when later post-training is incomplete or pushes in conflicting directions.

Although these are promising ideas, the available literature does not give a coherent picture of midtraining's effectiveness as a technique. For one thing, important design decisions such as data generation, data mixture curation, and training recipes vary widely and are not always available. For another, the reproducible literature has mostly trained small models on relatively small token budgets, making it difficult to understand how midtraining's effects scale with either parameter.

\paragraph{Our contributions.} We build a setting in which both hopes can be tested at once. In \emph{Dispatch}, models assign crews to trading runs. A seven-clause Charter and a profit rule usually pick the same crew, and we midtrain on a corpus that describes dispatchers applying the Charter. We then elicit the task with finetuning whose demonstrations are consistent with either motivation and exercise only five of the seven clauses (Figure~\ref{fig:robots_arrows}). We run this at scale, up to 190M midtraining tokens on a 110B-parameter model, with a 1B-token run in flight, and we report three results:

\begin{enumerate}
\item Asked what governs its decisions, the midtrained model says it follows the Charter, all seven clauses included.
\item It does not. On the five demonstrated clauses it follows the Charter 90\% of the time against 37\% for a control; on the two clauses the finetuning never exercised, 42\% against 10\% on GLM-4.5-Air and 15\% against 13\% on Gemma~3 27B. Midtraining installed the demonstrated clauses, not the Charter.
\item Even the demonstrated clauses depend on the finetuning data being perfectly consistent with the Charter. Replacing 2\% of the demonstrations with profit-labelled conflicts brings the rate to 13\%, and half a percent does most of the damage.
\end{enumerate}

We then vary everything else. More midtraining helps monotonically and has not saturated at 190M tokens; tenfold more elicitation data does not restore the held-out clauses; documents that describe the rule without a worked example install a much weaker prior; replacing supervised finetuning with \ac{rl} on the same episodes leaves the midtrained motivation retrievable but not load-bearing; and in a second setting, a fictional dialect of Python, held-out rules do transfer, which sharpens the question of what midtraining installs. We also report why reproductions of the published model-spec-midtraining recipe led us to build a new setting.

We read these results as a moderate downward update on alignment midtraining as currently practised in public, and as a specific one: the gap between what a midtrained model says and what it does is exactly the gap a developer cannot close by talking to the model.

% I want to say a few things here, which are all a bit fuzzy to me. First, we want to say that this is surprising! We should also give a bunch of caveats to the above results and not claim that they are necessarily representative of midtraining at the lab scales. Nonetheless, it's the closest that we have to our knowledge. Second, I want to say that these results are surprising and we want to understand them better! We are guessing we're doing something wrong compared to how it's done in labs, but it's not clear to us what this is! We should also say that doing similar deconflicting points seems valuable going forward.
